\chapter{2011年浙江卷（理）}

\section{单选题}

\begin{problem}
设函数 \( f(x) = \begin{cases} -x, & x \leqslant 0, \\ x^2, & x > 0 \end{cases} \) 若 \( f(\alpha) = 4 \)，则实数 \( \alpha = \)（\quad）

\choices
    {\(-4\) 或 \(-2\)}
    {\(-4\) 或 2}
    {\(-2\) 或 4}
    {\(-2\) 或 2}
\end{problem}

\begin{problem}
把复数 \(z\) 的共轭复数记作 \(\overline{z}\)，\(i\) 为虚数单位，若 \(z = 1 + \i\)，则 \((1 + z)\cdot \overline{z} =\)（\quad）

\choices
    {\(3 - \i\)}
    {\( 3 + \i \)}
    {\( 1 + 3\i \)}
    {\(3\)}
\end{problem}

\begin{problem}
若某几何体的三视图如图所示，则这个几何体的直观图可以是（\quad）

\begin{center}
\bitmapinclude[max width=0.88\linewidth,max totalheight=6.8cm]{img/2011/zhejiang_science/q3_options.png}
\end{center}
\end{problem}

\begin{problem}
下列命题中错误的是（\quad）

\choices
    {如果平面 \(\alpha \perp\) 平面 \(\beta\)，那么平面 \(\alpha\) 内一定存在直线平行于平面 \(\beta\)}
    {如果平面 \( \alpha \) 不垂直于平面 \( \beta \) ，那么平面 \( \alpha \) 内一定不存在直线垂直于平面 \( \beta \)}
    {如果平面 \(\alpha \perp\) 平面 \(\gamma\)，平面 \(\beta \perp\) 平面 \(\gamma\)，\(\alpha \cap \beta = l\)，那么 \(l \perp\) 平面 \(\gamma\)}
    {如果平面 \(\alpha \perp\) 平面 \(\beta\)，那么平面 \(\alpha\) 内所有直线都垂直于平面 \(\beta\)}
\end{problem}

\begin{problem}
设实数 \(x, y\) 满足不等式组 \(\left\{ \begin{array}{l} x + 2y - 5 > 0, \\ 2x + y - 7 > 0, \\ x \geqslant 0, y \geqslant 0. \end{array} \right.\) 若 \(x, y\) 为整数，则 \(3x + 4y\) 的最小值是（\quad）

\choices
    {14}
    {16}
    {17}
    {19}
\end{problem}

\begin{problem}
若 \(0 < \alpha < \frac{\pi}{2}, -\frac{\pi}{2} < \beta < 0, \cos \left(\frac{\pi}{4} + \alpha\right) = \frac{1}{3}, \cos \left(\frac{\pi}{4} - \frac{\beta}{2}\right) = \frac{\sqrt{3}}{3}\)，则 \(\cos \left(\alpha + \frac{\beta}{2}\right) =\)（\quad）

\choices
    {\(\frac{\sqrt{3}}{3}\)}
    {\(-\frac{\sqrt{3}}{3}\)}
    {\(\frac{5\sqrt{3}}{9}\)}
    {\(-\frac{\sqrt{6}}{9}\)}
\end{problem}

\begin{problem}
若 \(a, b\) 为实数，则 \(0 < ab < 1\) 是“\(a < \frac{1}{b}\) 或 \(b > \frac{1}{a}\)”的（\quad）

\choices
    {充分而不必要条件}
    {必要而不充分条件}
    {充分必要条件}
    {既不充分也不必要条件}
\end{problem}

\begin{problem}
已知椭圆 \(C_{1}:\frac{x^{2}}{a^{2}} +\frac{y^{2}}{b^{2}} = 1(a > b > 0)\) 与双曲线 \(C_2:x^2 -\frac{y^2}{4} = 1\) 有公共的焦点，\(C_2\) 的一条渐近线与以 \(C_1\) 的长轴为直径的圆相交于 \(A,B\) 两点.若 \(C_1\) 恰好将线段 \(AB\) 三等分，则（\quad）

\choices
    {\(a^2 = \frac{13}{2}\)}
    {\(a^2 = 13\)}
    {\(b^{2} = \frac{1}{2}\)}
    {\(b^{2} = 2\)}
\end{problem}

\begin{problem}
有 5 本不同的书，其中语文书 2 本，数学书 2 本，物理书 1 本. 若将其随机地抽取并摆放在书架的同一层上，则同一科目的书都不相邻的概率为（\quad）

\choices
    {\(\frac{1}{5}\)}
    {\(\frac{2}{5}\)}
    {\(\frac{3}{5}\)}
    {\(\frac{4}{5}\)}
\end{problem}

\begin{problem}
设 \(a, b, c\) 为实数，\(f(x) = (x + a)(x^2 + bx + c)\)，\(g(x) = (ax + 1)(cx^2 + bx + 1)\). 记集合 \(S = \{x \mid f(x) = 0, x \in \R\}\)，\(T = \{x \mid g(x) = 0, x \in \R\}\)，若 \(|S|\)，\(|T|\) 分别为集合 \(S, T\) 的元素个数，则下列结论不可能的是（\quad）

\choices
    {\( |S| = 1 \) 且 \( |T| = 0 \)}
    {\( |S| = 1 \) 且 \( |T| = 1 \)}
    {\( |S| = 2 \) 且 \( |T| = 2 \)}
    {\( |S| = 2 \) 且 \( |T| = 3 \)}
\end{problem}

\section{填空题}

\begin{problem}
若函数 \( f(x) = x^2 - |x + a| \) 为偶函数，则实数 \( a = \)  \fillinblank{}.
\end{problem}

\begin{problem}
某程序框图如图所示，则该程序运行后输出的 \(k\) 的值是 \fillinblank{}.

\begin{center}
\bitmapinclude[max width=0.28\linewidth,max totalheight=7.0cm]{img/2011/zhejiang_science/q12_fig1.png}
\end{center}
\end{problem}

\begin{problem}
设二项式 \(\left(x - \frac{a}{\sqrt{x}}\right)^6 (a > 0)\) 的展开式中 \(x^3\) 的系数为 \(A\)，常数项为 \(B\). 若 \(B = 4A\)，则 \(a\) 的值是 \fillinblank{}.
\end{problem}

\begin{problem}
若平面向量 \(\alpha, \beta\) 满足 \(|\alpha| = 1, |\beta| \leqslant 1\)，且以向量 \(\alpha, \beta\) 为邻边的平行四边形的面积为 \(\frac{1}{2}\)，则 \(\alpha\) 和 \(\beta\) 的夹角 \(\theta\) 的取值范围是 \fillinblank{}.
\end{problem}

\begin{problem}
某毕业生参加人才招聘会，分别向甲、乙、丙三个公司投递了个人简历，假定该毕业生得到甲公司面试的概率为 \(\frac{2}{3}\)，得到乙、丙公司面试的概率均为 \(p\)，且三个公司是否让其面试是相互独立的. 记 \(X\) 为该毕业生得到面试的公司个数. 若 \(P(X = 0) = \frac{1}{12}\)，则随机变量 \(X\) 的数学期望 \(E(X) = \fillinblank{}\).
\end{problem}

\begin{problem}
设 \(x, y\) 为实数. 若 \(4x^{2} + y^{2} + xy = 1\)，则 \(2x + y\) 的最大值是 \fillinblank{}.
\end{problem}

\begin{problem}
设 \(F_{1}, F_{2}\) 分别为椭圆 \(\frac{x^{2}}{3} + y^{2} = 1\) 的焦点，点 \(A, B\) 在椭圆上.若 \(\overrightarrow{F_{1}A} = 5\overrightarrow{F_{2}B}\)，则点 \(A\) 的坐标是 \fillinblank{}.
\end{problem}

\section{解答题}

\begin{problem}
在 \(\triangle ABC\) 中，角 \(A, B, C\) 所对的边分别为 \(a, b, c\). 已知 \(\sin A + \sin C = p \sin B (p \in \R)\)，且 \(ac = \frac{1}{4} b^2\). (1) 当 \(p = \frac{5}{4}, b = 1\) 时，求 \(a, c\) 的值; (2) 若角 B 为锐角，求 p 的取值范围.
\end{problem}

\begin{problem}
已知公差不为 0 的等差数列 \( \{a_{n}\} \) 的首项 \( a_{1} \) 为 \( a (a \in \R) \) . 设数列的前 n 项和为 \( S_{n} \) ，且 \( \frac{1}{a_{1}} , \frac{1}{a_{2}} , \frac{1}{a_{4}} \) 成等比数列. (1) 求数列 \( \{a_{n}\} \) 的通项公式及 \( S_{n} \) ; (2) 记 \(A_{n} = \frac{1}{S_{1}} + \frac{1}{S_{2}} + \frac{1}{S_{3}} + \cdots + \frac{1}{S_{n}}\)，\(B_{n} = \frac{1}{a_{1}} + \frac{1}{a_{2}} + \frac{1}{a_{2^{2}}} + \cdots + \frac{1}{a_{2^{n-1}}}\) 当 \(n \geqslant 2\) 时，试比较 \(A_{n}\) 与 \(B_{n}\) 的大小.
\end{problem}

\begin{problem}
如图，在三棱锥 \(P - ABC\) 中，\(AB = AC\)，\(D\) 为 \(BC\) 的中点，\(PO \perp\) 平面 \(ABC\)，垂足 \(O\) 落在线段 \(AD\) 上，已知 \(BC = 8\)，\(PO = 4\)，\(AO = 3\)，\(OD = 2\). (1) 证明: \(AP \perp BC\); (2) 在线段 \(AP\) 上是否存在点 \(M\)，使得二面角 \(A - MC - B\) 为直二面角. 若存在，求出 \(AM\) 的长; 若不存在，请说明理由.

\begin{center}
\bitmapinclude[max width=0.42\linewidth,max totalheight=6.6cm]{img/2011/zhejiang_science/q20_fig1.png}
\end{center}
\end{problem}

\begin{problem}
已知抛物线 \(C_1: x^2 = y\)，圆 \(C_2: x^2 + (y - 4)^2 = 1\) 的圆心为点 \(M\). (1) 求点 M 到抛物线 \( C_{1} \) 的准线的距离; (2) 已知点 \(P\) 是抛物线 \(C_{1}\) 上一点 (异于原点)，过点 \(P\) 作圆 \(C_{2}\) 的两条切线，交抛物线 \(C_{1}\) 于 \(A, B\) 两点，若过 \(M, P\) 两点的直线 \(l\) 垂直于 \(AB\)，求直线 \(l\) 的方程.

\begin{center}
\bitmapinclude[max width=0.42\linewidth,max totalheight=6.6cm]{img/2011/zhejiang_science/q21_fig1.png}
\end{center}
\end{problem}

\begin{problem}
设函数 \( f(x) = (x - a)^2 \ln x, a \in \R \). (1) 若 \(x = \e\) 为 \(y = f(x)\) 的极值点，求实数 \(a\); (2) 求实数 \(a\) 的取值范围，使得对任意的 \(x \in (0,3\e]\)，恒有 \(f(x) \leqslant 4\e^2\) 成立. 注：e 为自然对数的底数.
\end{problem}
